\chapter{How The DSL Works} \label{sec:HowTheDSLWorks}

\section{The World Representation} \label{sec:TheWorldRepresentation}

The world refers to the collection of data that defines the state of the characters and
their relations. It is split into \textbf{Actors}, which are either active contributors to your
story or passive props\slash rooms, and \textbf{Relations}, which include any sort of data attached
to or between these actors. One unique type of relation, which is explored more later,
is \textbf{Social Practices}, which provide the interface through which actors are able to interact
with each other.

\subsection{Actors}

Actors can take the shape of any entity within the world:

\begin{itemize}
    \item \textbf{Characters within the story:} Any character that acts autonomously and has
    individual desires or goals. This is the primary use of actors, and in a simple
    social scenario it may be all you need.
    \item \textbf{Static objects in the world:} For example, a house that characters can enter and
    exit, or a bottle that can be passed around. The simple example that comes pre-loaded
    in the playground and shown below uses this approach.
    \item \textbf{Games or other systems within the world:} Theoretically, one could even represent
    a dynamic chess board within the world as an actor with certain relation states inside
    of it. Or it could be a simpler case, like a jukebox that can play from a set number of
    songs. This language is by design very powerful; it is ease of use that is being evaluated.
    \item \textbf{Drama managers or other meta-entities:} One could introduce an all-powerful deistic
    entity that is able to create drama within character relations when things get too stale.
    This is commonly implemented within similarly structured projects as a “drama manager”,
    but here it would be just another actor within the world. All information is globally
    accessible, which enables this functionality.
\end{itemize}

Actors can be defined in the world script, along with a display name, goals, and modifiers,
as seen in \Cref{lst:ActorDef}. The display name is in quotes, and the lowercase specifier
is an internal ID that is used to refer to the actor throughout the simulation. Optional
modifiers and goals follow this, with goals being contained within brackets.

\begin{lstlisting}[float=hb, caption=Examples of instantiating agents into the world., label=lst:ActorDef]
actor house as "House" hidden
actor street as "Street" hidden

actor bob as "Bob" {
  goal (10): delta count Person where Person.feels.happy
}
actor alice as "Alice"
\end{lstlisting}

Any actors considered ``inactive'' will be marked as not being available to perform
actions. This is useful, for example, if a character becomes unconscious, but you still
want them to be displayed and other characters can still observe them. Actors marked
as ``hidden'' will still be able to perform actions (unless also marked as ``inactive''),
but will not be shown to the user. A drama manager might use this modifier to run its
operations in the background.

Actors can have goals that will influence their decision-making. Praxsmith uses a
planning system to determine which action an actor should take in order to maximize
the achievement of their goals, so goal definition is crucial in getting them to
behave correctly. The scores of different actions are available in the interface
for manual characters, and automatic characters will automatically pick the highest
scored option (ties broken randomly).

\Cref{lst:ActorDef} features one goal that tracks a numerical
value. Goals themselves are written as expressions, which are explored later in
this paper. The weight appears within parenthesis, so Bob's goal can be read as
``Bob wants to see the number of happy people increase with a weight of $10$''.
If this number of happy people changes, Bob's scoring of the resulting state
will reflect that, proportionally to the number of people in question. If two
people become happy, the resulting state will be $20$ points greater. If one
person stops being happy, the resulting state will be $-10$.

\subsection{Relations}

Relations are the primary method of tying actors together in the world. They can
represent actual human relations, such as “Alice trusts Bob”, but they can also
represent other concepts, like “Alice is in the house”.

\begin{lstlisting}[float=hb, caption=Example of defining and applying a new relation type., label=lst:RelationDef]
directional trusts {
  level: 0..10
}

alice.trusts.bob {
  level: 7
}
\end{lstlisting}

Definitions vary from type to type, but generally take the form of \Cref{lst:RelationDef}.
By following certain naming conventions, relations can be specified in a similar manner
to the English language. This example defines a non-mutual relation that allows for one actor to
trust another, and also gives it a ``level'' field that can range from 0 to 10. It then applies
the relation to Alice and Bob, stating that that Alice trusts Bob with a level of 7 out of 10.

One special kind of relation to note are practices, which represent social practices that actors
can engage in and have more rigid structures. They are the mechanism that gives characters the
ability to act and influence the world. More information about them is below.

\subsection{Sentences}

Sentences represent queries about the world state, and are used in expressions to pull
information from the world for use in various other mechanisms. They are heavily
derived from the sentences found in the original Praxis language this experiment is
inspired by. \cite{evans_ai_2015}

\begin{lstlisting}[float=hb, caption=Example of two sentences that query Alice's trust of Bob., label=lst:Sentence]
alice.trusts.bob
alice.trusts.bob.level
\end{lstlisting}

The first sentence acts as an existence check, returning true if the ``Alice trusts Bob''
relation exists at all. The second sentence takes that a step further, asking to what
level that trust exists in what is referred to as a ``fielded query''. Fielded queries
on relations that do not exist will throw an error, so it is important to first guard
against this case or the type checker will not allow the world to build.

\subsection{Expressions}

Expressions are the logical glue that allow for reasoning about the state of the world.
Sentences extract that information, and expressions allow you to ask more complex
questions about it. They take a form very similar to traditional mathematical equations
and predicate logic. Both Praxis \cite{evans_ai_2015} and Comme il Faut \cite{mccoy_comme_2011}
use predicate logic heavily as conditions, and Praxsmith is no different.

\begin{lstlisting}[float=hb, caption=Example expression that returns true if the number of happy or excited people is greater than 3., label=lst:Expression]
(count Person where (Person.feels.happy or Person.feels.excited)) > 3
\end{lstlisting}

In \Cref{lst:Expression}, the ``count'' operator is used to introduce a new binding, $Person$,
and the number of actors that match the following condition is compared in an inequality
with $3$. The introduction of free bindings like this makes expressions quite powerful,
although they do still come with limitations.

Numerous operators of different types are available to manipulate sentences and values
within expressions.

\begin{itemize}
    \item The $and$, $or$, and $not$ operators all behave according to traditional boolean
    logic. Praxsmith does not have any concept of ``truthy'' or ``falsy'' values; if the inputs
    are not booleans, the type checker will not allow the world to build. Importantly, both
    binary boolean operators short-circuit, which allows them to act as gates that guarantee
    the existence of certain relations before their fields are queried.
    \item The $is$ operator allows for two values to be compared for equality. Both sides
    must be the same type, or the type checker will fail.
    \item The $+$, $-$, $*$, and $/$ operators act primarily as numerical addition,
    subtraction, multiplication, and division, although inspiration was taken from
    modern scripting languages to allow for addition to also describe string conversion
    and appending. The multiplication also allows for more flexible operations, like
    string repetition and treating booleans as 1-or-0 values for a makeshift if-statement
    that returns a number.
    \item The $for all\ \_$ operator evaluates to true if and only if every actor, bound
    to the free variable, passes the condition. Care must be taken with this operator, as
    every actor also means the actor performing the action as well as static actors that
    would not count as characters, which must be filtered out if necessary.
    \item The $any\ \_\ where$ operator evaluates to true if any actor, bound to the free
    variable, passes the condition.
    \item The $count\ \_\ where$ evaluates to a number representing the number of actors,
    bound to the free variable, that pass the condition. This is one of the primary
    methods of moving from the boolean space to the number space, and allows for
    goals to more accurately track large numbers of actors.
    \item The $sum$, $average$, $min$, and $max$ operators represent aggregators that
    work similarly to the ``count'' operator, except instead of simply counting the number
    of actors that pass a condition, a certain value from those actors is accumulated
    according to the specified method. Degenerate cases default to $0$.
\end{itemize}

\subsection{Type Checking Guarantees}

When the world is built, all relevant statements are passed through the engine's type
checker to ensure conformity. This moves errors from runtime into build time, which is
crucial for a system with so many moving parts and possible outcomes.

\begin{lstlisting}[float=hb, caption=Examples describing how to appease the type checker., label=lst:TypeChecking]
// "cannot access field level of Other.trusts.Person without guarantee that it exists"
Other.trusts.Person.level > 5

// both compile
Other.trusts.Person and Other.trusts.Person.level > 5
not Other.trusts.Person or Other.trusts.Person.level < 5
\end{lstlisting}

\Cref{lst:TypeChecking} shows a couple of examples of expressions that either pass
or fail the type checking pass. Assuming that the $level$ field is in fact a number,
it may seem counterintuitive that the first expression fails. This is because it may
an error if the relation does not exist, which is a potential major pitfall,
especially when using the operators that iterate over large numbers of actors.

The other compiling examples both make use of the short-circuit guarantees of the $and$
and $or$ operators, in that the right-hand expression will only need to be evaluated
if the left-hand is $true$ or $false$, respectively. The type checker tracks these
guarantees along its iteration of the expression, thus ensuring that runtime errors are
never thrown due to invalid field accesses.

This also allows for custom default/fallback behavior. The first compiling example in
\Cref{lst:TypeChecking} requires that $Other$ trusts $Person$ with a $level$ of greater than
$5$, while the second compiling example requires that $Other$ trusts $Person$ with a $level$
of less than $5$ or, importantly, not at all. This behavior is important to keep in mind
when writing expressions, and the condensed nature of this predicate-style logic can be hard
to keep track of when first using the language.

\section{Relation Types} \label{sec:RelationTypes}

Relations in Praxsmith come in many different types, and each behaves very differently
from the others. They are mainly derived from the emergent types the authors of Praxis
noticed \cite{evans_ai_2015}, and many of the names are derived from the social fact
types of the Ensemble Engine \cite{samuel_ensemble_2015}.

\begin{lstlisting}[float=hb, caption=Examples of how to instantiate the simple relation types., label=lst:RelationTypeExamples]
// Defining traits
trait curious
// Applying / querying traits
bob.is.curious

// Defining emotions
emotion happy
// Applying / querying emotions
alice.feels.happy

// Defining directionals
directional trusts
exclusive directional is_in
// Applying / querying directionals
alice.trusts.bob
Actor.is_in.house

// Defining reciprocals
reciprocal is_married_to
// Applying / querying reciprocals (both the following are equivalent)
alice.is_married_to.bob
bob.is_married_to.alice
\end{lstlisting}

\textbf{Traits} define a property that an actor can have. They only ever have one connected
actor, which is the actor expressing that trait. An actor can have many traits,
and they can be added or removed at any time. As seen in \Cref{lst:RelationTypeExamples},
they can be accessed using the $is$ keyword.

\textbf{Emotions} define a short-lived feeling an actor can have. Similarly to traits, they
only ever have one connected actor, which is the actor expressing that emotion.
Unlike traits, an actor can only ever have one emotion at a time. If long-term effects
are desired, this is not the relation type to use. Emotions are intended as a quick
way to represent how a character is feeling in the moment, and can be easily queried
and displayed by the frontend for the player to see. As seen in \Cref{lst:RelationTypeExamples},
they can be accessed using the $feels$ keyword.

\textbf{Directional relations} connect two actors asymmetrically. In the first
directional application example in \Cref{lst:RelationTypeExamples}, Alice
trusts Bob, but that does not necessarily mean that Bob trusts Alice. Most relations
present this way, in which one character feels a certain way about another, but that
feeling is not necessarily mutual.

By adding exclusive to the type definition, the engine will ensure that the primary actor
(in the case of \Cref{lst:RelationTypeExamples}, Alice) only has one outgoing
relation of this type at any given time. For example, Alice can only be in one
place at a time, so it doesn't make sense for her to have more than one $is\_in$ relation.
If a new relation is created that conflicts, it will automatically remove the old one.
This is very useful for switching the state of a character between mutually exclusive states,
such as them moving from one room to another. Without this, one would have to manually
remove their old location relation every time they moved, which would get very
tedious and is prone to error.

Sometimes, relationships are inherently mutual. For example, if Alice is married to Bob,
then Bob is married to Alice. \textbf{Reciprocal} relations are designed for this case, and the
engine will ensure that if and only if one edge exists, the other does as well. Any
fields attached to a reciprocal relation are shared between the two edges, as emphasized
on the last couple lines of \Cref{lst:RelationTypeExamples}.

\subsection{Practices}

Practices are a special type of relationship that give your characters to perform actions.

\begin{lstlisting}[float=hb, caption=A simple example practice allowing for two characters to greet each other., label=lst:ExamplePractice]
emotion happy

exclusive directional is_in

practice greet (Greeter, Greeted) {
  exclamation: string
  actions: [
    {
      for: Greeter
      name: "Greet [Greeted]"
      conditions: [
        self
        any Place where Greeter.is_in.Place and Greeted.is_in.Place
      ]
      outcomes: [
        say self.exclamation
        set Greeted.feels.happy
        broadcast "[Greeter]'s greeting has made [Greeted] happy!"
        delete self
      ]
    }
  ]
}

// must also be instantiated on the world side:
practice.greet.bob.alice { exclamation: "Hey [Greeted]!" }
\end{lstlisting}

All practices have participants. In \Cref{lst:ExamplePractice}, there are two: $Greeter$
and $Greeted$, and each will play a unique role in the practice. This behavior is
common across all practices. These parameters ($Greeter$ and $Greeted$) will then become
bound to $bob$ and $alice$ respectively for any expressions or strings within that
particular instance, as can be seen throughout the example, where $Greeter$ and
$Greeted$ are treated as regular actors in the expression despite not being defined
in the world. The example also specifies an $exclamation$ field, which
allows for the specific greeting to be customized. This is another use of fields,
alongside tracking changing state, which allows for characters to be individually
expressive even when their actions are driven by generic practices.

As seen in \Cref{lst:ExamplePractice}, practices also create a $self$ binding for any
expressions within them, allowing for easy access to internal fields or easy self-removal
after practice completion.

Each practice provides a set of actions that allow actors involved in that practice
to act. Each action will only be available if all of its conditions evaluate to $true$,
as dictated by the mandatory $conditions$ field.

There are four total mandatory fields in actions:

\begin{itemize}
    \item The $for$ field specifies which participant the action should be available
    for; in this case, greeting is an action available to $Greeter$. An action can
    only be made available to one participant, so if you want another similar action,
    it must be written explicitly with the participants flipped. This could be seen
    as a limitation, but it also helps ensure that actor-dependent outcomes don't get
    mixed up in a single action and is rarely an issue in practice.
    \item The $name$ field represents the display name of the action, as it will
    appear to the player. The display names of parameters can be inserted here by
    putting them in square brackets, as shown in the example.
    \item The $conditions$ field is a list of expressions that must all evaluate to
    true for the action to be available. Any relation existences guaranteed here will
    carry over as guarantees for outcomes.
    \item The $outcomes$ field contains a series of world state changes that occur when
    this action is performed. Options for mutations are detailed below.
\end{itemize}

\Cref{lst:ExamplePractice} only has one action listed, but practices may have as many as
they would like, and they often do in more complex scenarios.

Refered to internally as effects, these are the possible outcomes an action can have:

\begin{itemize}
    \item The $broadcast$ and $say$ effects allow for text to be printed to the screen.
    If the text is broadcasted, it will appear as system text, otherwise it will appear
    as though the character themselves is saying the line. The line can be any expression
    that evaluates to a string, which allows for fairly expressive speech, as explored
    in \Cref{lst:ExamplePractice}.
    \item The $activate$ and $deactivate$ effects flag actors as either able or unable
    to perform actions. For example, this might be used to stop an unconscious person
    from taking further action.
    \item The $hide$ and $unhide$ effects flag actors as either visible or invisible
    on the frontend. This may be used to reveal the existence of a previously unknown
    character.
    \item The $delete$ effect removes a relation from the world. Actors cannot be deleted.
    This is commonly used as $delete self$ to remove practices after they have served
    their course.
    \item The $create$ effect instantiates a new relation into the world. If an identical
    relation already exists, it will be replaced with this new one. This is how relations
    create other relations, which is an important mechanism in keeping the world running.
    For example, if an argue practice escalates to a high degree, an action may become
    available for one of the participants that allows for them to initiate a divorce
    practice if they are married.
    \item The $set$ effect sets a single field of a relation to a new value, which can be
    any expression, as long as it resolves to something within the target field type.
    \item The $increase$ and $cycle$ effects are both used to increase or decrease the
    value of a number or state of a variant, and behave differently on the boundaries of
    the range. Cycle will cause numbers to wrap around the ends of their ranges and variants
    to jump to the other end once exhausted, while increase will stop at boundaries.
\end{itemize}

\subsection{Field Types}

There are five types a field can be:

\begin{itemize}
    \item Numeric ranges: A field can be defined as a numeric range, such as $0..10$.
    This means that the field can take on any value between $0$ and $10$, inclusive.
    Numbers are stored as decimal values.
    \item Variant lists: A field can also be defined as a list of named variants,
    such as $calm\ \text{\textbar}\ tense\ \text{\textbar}\ furious$. This means that the field can only take on one
    of those three values. This kind of value is often referred to as an enum in other
    languages. When assigning to or accessing this field, to avoid collisions with
    sentences, the language requires you prefix the value with a $\#$.
    \item Actor References: A field can reference an actor in the world, and will
    be type checked upon world load to make sure the actor exists or the binding is
    valid. When assigning to or accessing this field, to avoid collisions with
    sentences, the language requires you to prefix the value with a $@$.
    \item String: A field can be a string, which can take on any value. When assigning
    to or accessing this field, the value must be wrapped in quotes.
    \item Boolean: A field can be a boolean, which can be either $true$ or $false$.
\end{itemize}

\section{Simple Example} \label{sec:SimpleExample}

\Cref{lst:SimpleExampleTypes} and \Cref{lst:SimpleExampleWorld} are meant to showcase a few core parts of the language in
a condensed manner. It includes two characters, Bob and Alice. Bob wants to make
people happy, but can only do so by greeting Alice, which he can only do if they are
in the same place. Alice starts off feeling sad and is on the street, while Bob is
in the house. To end the simulation, Alice must first move to the house (which is
permitted via the $practice.move.alice.street.house$ declaration that puts her in that
practice at initialization), and then Bob can greet her to make her happy, which
satisfies one of his goals.

\begin{lstlisting}[float=hb, caption=Type definitions for the simple example., label=lst:SimpleExampleTypes]
emotion sad
emotion happy

exclusive directional is_in

practice greet (Greeter, Greeted) {
  exclamation: string
  actions: [
    {
      for: Greeter
      name: "Greet [Greeted]"
      conditions: [
        self
        any Place where Greeter.is_in.Place and Greeted.is_in.Place
      ]
      outcomes: [
        say self.exclamation
        set Greeted.feels.happy
        broadcast "[Greeter]'s greeting has made [Greeted] happy!"
        delete self
      ]
    }
  ]
}

practice move (Actor, From, To) {
  actions: [
    {
      for: Actor
      name: "Move from [From] to [To]"
      conditions: [
        Actor.is_in.From
      ]
      outcomes: [
        set Actor.is_in.To
        broadcast "[Actor] moves from [From] to [To]."
      ]
    }
  ]
}
\end{lstlisting}

\begin{lstlisting}[float=hb, caption=World initialization for the simple example., label=lst:SimpleExampleWorld]
actor house as "House" hidden
actor street as "Street" hidden

actor bob as "Bob" {
  goal (10): delta count Person where Person.feels.happy
}
actor alice as "Alice"

bob.is_in.house
alice.is_in.street

alice.feels.sad

practice.greet.bob.alice {
  exclamation: "Hey [Greeted]! Welcome home!"
}
practice.move.alice.street.house
\end{lstlisting}

\section{The User Interface} \label{sec:TheUserInterface}

Alongside the engine, a simple webpage was created to facilitate interaction with
Praxsmith worlds, as well as their development. This included an authoring page where
syntax-highlighted code windows were supplied for both the world and type scripts,
as well as a gameplay page where experiences could be played. Actors could be set
to either manual or automatic, where manual actors could be controlled by the player
and automatic actors would pick the highest scoring action, with ties broken randomly.
A chat log was provided for anything said or broadcast during gameplay, and actor
emotions were displayed next to their names if available. Detailed error messages were
displayed for any issues encountered during any stage of the experience, and a debug
panel was provided for real-time queries and world mutations. Screenshots of both pages
to the interactive part of the project can be found in \Cref{fig:Webpage1} and \Cref{fig:Webpage2}.

\begin{figure}
    \centering
    \includegraphics[width=6in]{praxsmith_webpage_1}
    \caption{The authoring interface of the website.}
    \label{fig:Webpage1}
\end{figure}

\begin{figure}
    \centering
    \includegraphics[width=6in]{praxsmith_webpage_2}
    \caption{The gameplay window of the website.}
    \label{fig:Webpage2}
\end{figure}

There was also significant effort put into the documentation, which is linked on the
main page. This comprises of detailed explanations of all that is discussed in this
chapter of this research paper, as well as more examples to help new users get started.

As long as it remains online, the web interface for this project can be found at
https://praxsmith.emuman.net/.
